\chapter{Introducción}

\section{Motivación y Contexto}
En los últimos años, la sociedad ha sido testigo de un aumento considerable en el uso de la Inteligencia Artificial (IA) para suplantar la identidad de personas de forma maliciosa mediante la clonación de voz, técnica conocida como \textit{deepfake}. Esta evolución tecnológica ha alcanzado un punto crítico donde la distinción entre un audio humano y uno sintético es prácticamente imperceptible para el oído humano. Como se cuenta en ~\parencite{Barrington2025, Nadine2025}, estudios recientes indican que, mientras la detección de audios 100\% generados (i.e. voz puramente sintética, sin clonar otra voz) resulta sencilla para un usuario promedio, la precisión al clasificar voces reales frente a clones de alta calidad apenas ronda el 60\%.

Esta vulnerabilidad supone una amenaza latente en sectores estratégicos como la ciberseguridad, el periodismo y el ámbito judicial. Por otro lado, las técnicas de detección de \textit{deepfakes} no han avanzado al mismo ritmo que las de generación, lo que se traduce en una proliferación de estafas y desinformación. Por este motivo, ya no es posible confiar únicamente en las capacidades naturales del ser humano para diferenciar la autenticidad de una voz, lo que convierte el desarrollo de herramientas tecnológicas robustas en una necesidad social imperativa.

Especialmente ha crecido la importancia de diferenciar entre una voz real y una falsa en las ciencias forenses, ya que es vital, ante un hecho delictivo, el poder identificar correctamente a las distintas partes involucradas. Esto se explica en ~\parencite{Peña2024}, donde se relaciona la fonética forense y la criminalística para, entre otras cosas, poder decir si dos voces pertenecen a la misma persona. Este problema es muy importante ya que la voz de una misma persona puede verse alterada por patologías, ya sea de forma temporal o permanente, o variar dependiendo del idioma en el que esté hablando. También entran en juego factores socioeconómicos y hábitos alimenticios. Todo esto influye en la voz, y es crucial tenerlo en cuenta en las investigaciones para poder señalar al culpable.

\section{Definición del Problema}
El rendimiento del modelo de clasificación no es el único factor determinante en esta investigación. En la actualidad, los modelos de aprendizaje profundo se han convertido en ``cajas negras'', sistemas opacos donde no se comprende cómo se generan las predicciones. Esta falta de transparencia da lugar a problemas de confianza respecto a dichas decisiones, especialmente en ámbitos críticos como la medicina o la ingeniería, donde la vida de las personas puede estar en juego.

Para solventar este problema ha surgido la IA explicable (XAI), la cual toma dichos modelos ``caja negra'' y proporciona una justificación de qué características de los datos afectan más a la decisión del modelo. En esta investigación, abordamos el problema dual de maximizar la capacidad de discriminación entre audios reales y generados y, simultáneamente, dotar al sistema de una interpretación humana coherente mediante la aplicación de XAI.

\section{Objetivos}
El objetivo principal de este trabajo es desarrollar un marco de trabajo para la detección de audio generado sintéticamente, que priorice no solo la eficacia en la clasificación, sino también la transparencia de sus resultados. Para ello, se plantean los siguientes objetivos específicos:

\begin{itemize}
    \item \textbf{Exploración de Modelos de Clasificación:} Evaluar la viabilidad de distintos enfoques algorítmicos, desde métodos estadísticos tradicionales hasta arquitecturas de aprendizaje profundo de última generación, para identificar cuál se adapta mejor a la naturaleza del audio clonado.
    \item \textbf{Estudio de Características Acústicas:} Investigar qué rasgos del habla (frecuencia, timbre, periodicidad o pausas) contienen la información más relevante para distinguir la huella de una IA frente a la de un humano.
    \item \textbf{Integración de Métodos de Explicabilidad:} Proponer el uso de herramientas de interpretación que permitan visualizar qué partes del audio están influyendo en la decisión del sistema, convirtiendo el modelo en una herramienta útil para un análisis humano posterior.
\end{itemize}

\section{Planificación}
Para garantizar la consecución de los objetivos propuestos, el proyecto se ha estructurado siguiendo un flujo de trabajo sistemático. La metodología planificada se divide en las siguientes fases consecutivas:

\begin{enumerate}
    \item \textbf{Investigación de modelos y técnicas:} Estudio exhaustivo del estado del arte para identificar las arquitecturas de clasificación y los métodos de explicabilidad más prometedores en el dominio del procesamiento de voz y la detección de anomalías.
    \item \textbf{Adquisición y preparación de los datos:} Selección y procesamiento de un conjunto de datos equilibrado que contenga muestras de audio tanto reales como sintéticas, asegurando la correcta extracción de características y la adecuación de los formatos para el entrenamiento.
    \item \textbf{Prueba de los modelos y técnicas:} Fase de experimentación preliminar donde se evalúa el comportamiento individual de distintos algoritmos, ajustando parámetros básicos y analizando su capacidad inicial para discriminar entre las clases propuestas.
    \item \textbf{Implementación del pipeline completo:} Integración de los componentes de preprocesamiento, clasificación y explicabilidad en un sistema unificado, permitiendo una transición fluida desde la entrada del audio en crudo hasta la generación de un veredicto justificado.
    \item \textbf{Análisis de los resultados:} Evaluación crítica del sistema mediante métricas de rendimiento y auditorías de interpretabilidad, comparando los hallazgos obtenidos con las hipótesis iniciales para extraer conclusiones sólidas sobre la viabilidad del enfoque propuesto.
\end{enumerate}
